\documentclass[aspectratio=169,11pt]{beamer}

% ============================================================
% PACOTES
% ============================================================

\usepackage[utf8]{inputenc}
\usepackage[T1]{fontenc}
\usepackage[brazil]{babel}

\usepackage{amsmath}
\usepackage{amssymb}
\usepackage{graphicx}
\usepackage{tikz}
\usepackage[most]{tcolorbox}
\usepackage{lmodern}
\usepackage{array}
\usepackage{tabularx}
\usepackage{booktabs}

\usetikzlibrary{
    shadows,
    shadows.blur
}

% ============================================================
% CORES
% ============================================================

\definecolor{mainblue}{HTML}{3B35B9}
\definecolor{darkblue}{HTML}{16175A}

\definecolor{lightbox}{HTML}{ECECF7}

\definecolor{alertred}{HTML}{D60000}
\definecolor{alertbg}{HTML}{F8E5E5}

\definecolor{softgray}{HTML}{F7F7F9}

% ============================================================
% CONFIGURAÇÃO GERAL DO BEAMER
% ============================================================

\setbeamertemplate{navigation symbols}{}

\setbeamercolor{background canvas}{bg=white}
\setbeamercolor{normal text}{fg=black}
\setbeamercolor{item}{fg=mainblue}

\setbeamerfont{frametitle}{size=\Large,series=\mdseries}
\setbeamerfont{normal text}{size=\normalsize}

\setbeamersize{text margin left=0.60cm,text margin right=0.60cm}

\setbeamertemplate{itemize item}{\color{mainblue}$\blacktriangleright$}
\setbeamertemplate{itemize subitem}{\color{mainblue}$\blacktriangleright$}

\setlength{\leftmargini}{1.5em}
\setlength{\leftmarginii}{1.3em}

% ============================================================
% FAIXA SUPERIOR / TÍTULO DO SLIDE
% ============================================================

\setbeamercolor{frametitle}{bg=mainblue,fg=white}

\setbeamertemplate{frametitle}{%
    \nointerlineskip%
    \begin{beamercolorbox}[
        wd=\paperwidth,
        ht=0.72cm,
        dp=0.17cm,
        leftskip=0.32cm
    ]{frametitle}%
        \usebeamerfont{frametitle}%
        \insertframetitle%
    \end{beamercolorbox}%
}

% ============================================================
% RODAPÉ
% ============================================================

\setbeamercolor{footerleft}{bg=darkblue,fg=white}
\setbeamercolor{footerright}{bg=mainblue,fg=white}

\newcommand{\footersafespace}{0.30cm}
\newcommand{\footerheight}{0.25cm}
\newcommand{\footerdepth}{0.09cm}

\setbeamertemplate{footline}{%
    \leavevmode%
    \vbox{%
        \vskip\footersafespace%
        \hbox to \paperwidth{%
            \begin{beamercolorbox}[
                wd=0.85\paperwidth,
                ht=\footerheight,
                dp=\footerdepth,
                leftskip=0.10cm
            ]{footerleft}%
                \tiny
                INF01146 -- Avaliação de Desempenho -- Trabalho Final: Etapa 1%
            \end{beamercolorbox}%
            \begin{beamercolorbox}[
                wd=0.15\paperwidth,
                ht=\footerheight,
                dp=\footerdepth,
                rightskip=0.10cm
            ]{footerright}%
                \hfill
                \tiny
                \insertframenumber/\inserttotalframenumber%
            \end{beamercolorbox}%
        }%
    }%
}

% ============================================================
% CAIXAS
% ============================================================

\tcbset{
    slidebox/.style={
        enhanced,
        boxrule=0pt,
        arc=2mm,
        outer arc=2mm,
        left=2.2mm,
        right=2.2mm,
        top=1.4mm,
        bottom=1.4mm,
        toptitle=1mm,
        bottomtitle=1mm,
        drop shadow={
            shadow xshift=1.4mm,
            shadow yshift=-1.4mm,
            opacity=0.35
        }
    }
}

\newtcolorbox{questionbox}[1]{
    slidebox,
    colback=lightbox,
    colframe=lightbox,
    coltitle=white,
    colbacktitle=darkblue,
    title=#1,
    fonttitle=\normalsize
}

\newtcolorbox{impactbox}[1]{
    slidebox,
    colback=alertbg,
    colframe=alertbg,
    coltitle=white,
    colbacktitle=alertred,
    title=#1,
    fonttitle=\normalsize
}

\newcommand{\sectionframe}[1]{%
    {%
    \setbeamertemplate{frametitle}{}%
    \begin{frame}
        \centering
        \vfill
        {\Large #1}
        \vfill
    \end{frame}
    }%
}

\newcommand{\bits}[1]{\texttt{#1}}

% ============================================================
% METADADOS
% ============================================================

\title{Avaliação de Desempenho Computacional}
\subtitle{Algoritmo Genético -- Etapa 1}

\author{
    Felipe Avencourt Soares -- 342904\\
    Leonardo Rocha dos Santos -- 341831
}
\institute{Instituto de Informática -- UFRGS}
\date{2026/2}

% ============================================================
% DOCUMENTO
% ============================================================

\begin{document}

% ============================================================
% CAPA
% ============================================================

{%
\setbeamertemplate{frametitle}{}%

\begin{frame}

    \vfill

    \begin{center}

        \begin{tcolorbox}[
            enhanced,
            width=0.92\textwidth,
            colback=mainblue,
            colframe=mainblue,
            arc=2mm,
            boxrule=0pt,
            left=3mm,
            right=3mm,
            top=5mm,
            bottom=5mm,
            drop shadow={
                shadow xshift=1.7mm,
                shadow yshift=-1.7mm,
                opacity=0.35
            }
        ]

            \centering

            {\color{white}\LARGE
                \inserttitle
                \par
            }

            \vspace{0.15cm}

            {\color{white}\Large
                \insertsubtitle
            }

        \end{tcolorbox}

        \vspace{0.75cm}

        {\large
            \insertauthor
        }

        \vspace{0.35cm}

        {\small
            \insertinstitute
        }

        \vspace{0.35cm}

        {\small
            \insertdate
        }

    \end{center}

    \vfill

\end{frame}
}

% ============================================================
% SLIDE 2
% ============================================================

\begin{frame}{Ideia central}

    \vspace{0.15cm}

    \begin{questionbox}{Proposta}
        Caracterizar o desempenho computacional de um
        Algoritmo Genético sequencial ao variar o tamanho da
        população e o tamanho dos cromossomos.
    \end{questionbox}

    \vspace{0.28cm}

    \textbf{Não iremos:}

    \begin{itemize}
        \item provar que o AG encontra soluções melhores;
        \item comparar versões sofisticadas de Algoritmos Genéticos;
    \end{itemize}

    \vspace{0.10cm}

    \begin{impactbox}{Foco do trabalho}
        O AG se trata apenas de uma carga de trabalho computacional controlável.
    \end{impactbox}

\end{frame}

% ============================================================
% SLIDE 3
% ============================================================

\begin{frame}{O que é um Algoritmo Genético?}

    Um \textbf{Algoritmo Genético (AG)} é uma técnica de busca que mantém várias soluções
    candidatas ao mesmo tempo.

    \vspace{0.28cm}

    \begin{columns}[T,onlytextwidth]

        \begin{column}{0.48\textwidth}

            \vspace{-0.18cm}

            \begin{questionbox}{Ideia intuitiva}
                Em vez de testar uma solução por vez, o algoritmo mantém uma população de
                soluções e tenta melhorá-la geração após geração.
            \end{questionbox}

        \end{column}

        \begin{column}{0.46\textwidth}

            \textbf{Lógica geral:}

            \begin{itemize}
                \item avalia as soluções atuais;
                \item escolhe as melhores como base;
                \item combina e altera essas soluções;
                \item repete o processo.
            \end{itemize}

        \end{column}

    \end{columns}

    \vspace{0.18cm}

    \begin{impactbox}{Importante}
        O AG se baseia em vetores, números aleatórios e laços de repetição. São apenas
        operações computacionais inspiradas na evolução biológica.
    \end{impactbox}

\end{frame}

% ============================================================
% SLIDE 4
% ============================================================

\begin{frame}{Conceitos principais de um AG}

    \vspace{0.05cm}

    \renewcommand{\arraystretch}{1.25}
    \begin{tabularx}{\textwidth}{>{\bfseries}p{0.27\textwidth} X}
        Indivíduo & Uma solução candidata para o problema. \\
        \hline
        Cromossomo & A forma como o indivíduo é armazenado no programa. Neste trabalho, será um vetor binário. \\
        \hline
        População & Conjunto de $N$ indivíduos mantidos simultaneamente. \\
        \hline
        Fitness & Valor numérico que mede a qualidade de um indivíduo. \\
        \hline
        Geração & Uma iteração completa: avaliar, selecionar, cruzar, mutar e criar nova população. \\
    \end{tabularx}

    \vspace{0.25cm}

    \begin{questionbox}{Exemplo}
        O cromossomo \bits{101101} representa um indivíduo binário com $L=6$.
    \end{questionbox}

\end{frame}

% ============================================================
% SLIDE 5
% ============================================================

\begin{frame}{Representação usada no Trabalho}

    \vspace{0.05cm}

    Cada indivíduo será um vetor de bits com comprimento $L$.

    \vspace{0.18cm}

    \[
        \underbrace{[1,0,1,1,0,1]}_{\text{um indivíduo}}
        \qquad L=6
    \]

    \vspace{0.15cm}

    A população contém $N$ indivíduos:

    \vspace{0.05cm}

    \[
        \begin{array}{c}
        \bits{101101} \\
        \bits{001011} \\
        \bits{111000} \\
        \vdots \\
        \bits{010110}
        \end{array}
        \qquad
        \left\}
        \begin{array}{l}
        N \text{ linhas} \\
        L \text{ bits por linha}
        \end{array}
        \right.
    \]

    \vspace{0.08cm}

    \textbf{Relação com desempenho:}
    aumentar $N$ ou $L$ aumenta dados manipulados e trabalho por geração.

\end{frame}

% ============================================================
% SLIDE 6
% ============================================================

\begin{frame}{Fitness OneMax: cálculo da aptidão}

    \vspace{0.10cm}

    A função de fitness escolhida será \textbf{OneMax}.

    \vspace{0.18cm}

    Ela conta quantos bits iguais a 1 existem no cromossomo.

    \vspace{0.25cm}

    \[
        x = [1,0,1,1,0,1]
    \]

    \[
        fitness(x) = 1 + 0 + 1 + 1 + 0 + 1 = 4
    \]

    \vspace{0.15cm}

    \begin{questionbox}{Por que usar OneMax?}
        É simples, controlada e fácil de verificar. Assim, o problema de
        otimização não esconde o comportamento computacional do AG.
    \end{questionbox}

\end{frame}

% ============================================================
% SLIDE 7
% ============================================================

\begin{frame}{Etapas de uma geração}

    \vspace{0.15cm}

    Em cada geração, o AG executa as seguintes etapas:

    \vspace{0.15cm}

    \begin{enumerate}
        \item \textbf{Avaliação}: calcula o fitness de cada indivíduo.
        \item \textbf{Seleção}: escolhe indivíduos para serem pais.
        \item \textbf{Crossover}: combina partes dos pais para gerar filhos.
        \item \textbf{Mutação}: altera alguns bits dos filhos.
        \item \textbf{Substituição}: a nova população passa a ser usada na próxima geração.
    \end{enumerate}

    \vspace{0.20cm}

    \[
        \text{População atual}
        \rightarrow
        \text{avaliação}
        \rightarrow
        \text{seleção}
        \rightarrow
        \text{crossover}
        \rightarrow
        \text{mutação}
        \rightarrow
        \text{nova população}
    \]

    \begin{impactbox}{No experimento}
        O número de gerações $G$ será fixo para que todas as configurações
        executem uma quantidade comparável de trabalho.
    \end{impactbox}

\end{frame}

% ============================================================
% SLIDE 8
% ============================================================

\begin{frame}{Seleção por torneio}

    \vspace{0.10cm}

    A seleção define quais indivíduos serão usados como pais.

    \vspace{0.20cm}

    \begin{questionbox}{Torneio de tamanho 2}
        Escolhem-se dois indivíduos aleatoriamente. O que tiver maior fitness
        vence e é selecionado como pai.
    \end{questionbox}

    \vspace{0.25cm}

    \[
        \begin{array}{rcl}
        A &=& \bits{101101} \quad fitness(A)=4 \\
        B &=& \bits{001011} \quad fitness(B)=3
        \end{array}
    \]

    \vspace{0.10cm}

    \[
        fitness(A) > fitness(B) \quad \Rightarrow \quad A \text{ vence o torneio}
    \]

    \vspace{0.08cm}

    \textbf{Efeito computacional:}
    acessa valores de fitness e realiza comparações.

\end{frame}

% ============================================================
% SLIDE 9
% ============================================================

\begin{frame}{Crossover: combinando dois pais}

    \vspace{0.10cm}

    O crossover gera novos indivíduos combinando partes de dois pais.

    \vspace{0.20cm}

    \begin{questionbox}{Crossover de um ponto}
        Escolhe-se uma posição de corte. O filho recebe a primeira parte de
        um pai e a segunda parte do outro.
    \end{questionbox}

    \vspace{0.25cm}

    \[
        \begin{array}{rcl}
        Pai\;1 &=& \bits{101\,|\,101} \\
        Pai\;2 &=& \bits{001\,|\,011} \\
        \hline
        Filho &=& \bits{101\,|\,011}
        \end{array}
    \]

    \vspace{0.08cm}

    \textbf{Efeito computacional:}
    copia e combina posições do cromossomo; custo cresce com $L$.

\end{frame}

% ============================================================
% SLIDE 10
% ============================================================

\begin{frame}{Mutação: pequenas alterações aleatórias}

    \vspace{0.10cm}

    A mutação altera alguns bits dos filhos após o crossover.

    \vspace{0.20cm}

    \begin{questionbox}{Mutação bit-flip}
        Cada posição pode ser invertida: $0 \rightarrow 1$ ou $1 \rightarrow 0$.
    \end{questionbox}

    \vspace{0.25cm}

    \[
        \bits{101011}
        \quad \xrightarrow{\text{inverte a posição 4}} \quad
        \bits{101111}
    \]

    \vspace{0.20cm}

    \begin{impactbox}{Efeito computacional}
        A mutação envolve percorrer cromossomos e gerar decisões aleatórias.
        Mesmo sendo conceitualmente simples, pode ter custo relevante.
    \end{impactbox}

\end{frame}

% ============================================================
% SLIDE 11
% ============================================================

\begin{frame}{Lógica computacional do AG}

    \vspace{0.05cm}

    Para uma população com $N$ indivíduos e cromossomos de tamanho $L$:

    \vspace{0.15cm}

    \begin{itemize}
        \item a avaliação percorre aproximadamente $N \times L$ bits por geração;
        \item crossover e mutação também manipulam cromossomos de tamanho $L$;
        \item com $G$ gerações, a carga nominal cresce com $N \times L \times G$.
    \end{itemize}

    \vspace{0.12cm}

    \[
        \text{carga nominal} \approx N \times L \times G
    \]

    \vspace{0.02cm}

    \begin{questionbox}{Exemplo}
        Se $N=1000$, $L=256$ e $G=100$, então
        $NLG = 25.600.000$ posições de cromossomo são consideradas
        na normalização.
    \end{questionbox}

\end{frame}

% ============================================================
% SLIDE 12
% ============================================================

\begin{frame}{Objeto computacional}

    \begin{columns}[T,onlytextwidth]

        \begin{column}{0.49\textwidth}

            \textbf{AG geracional sequencial}

            \begin{itemize}
                \item Implementação em C++
                \item Execução em Linux
                \item Uma única thread
                \item Número fixo de gerações $G$
            \end{itemize}

        \end{column}

        \begin{column}{0.48\textwidth}

            \textbf{Configuração simplificada}

            \begin{itemize}
                \item Fitness OneMax
                \item Seleção por torneio
                \item Crossover de um ponto
                \item Mutação bit-flip
            \end{itemize}

        \end{column}

    \end{columns}

    \vspace{0.25cm}

    \begin{impactbox}{Delimitação}
        A complexidade do trabalho está na avaliação de desempenho computacional,
        não na sofisticação do Algoritmo Genético.
    \end{impactbox}

\end{frame}

% ============================================================
% SLIDE 13
% ============================================================

\begin{frame}{Questão de pesquisa}

    \vspace{0.15cm}

    \begin{questionbox}{Pergunta principal}
        Como o aumento de $N$ e $L$ afeta o desempenho computacional
        de um Algoritmo Genético sequencial?
    \end{questionbox}

    \vspace{0.28cm}

    \textbf{Pergunta operacional:}

    \begin{itemize}
        \item O custo cresce proporcionalmente a $N \times L$?
        \item Ou surgem regimes diferentes quando o conjunto de trabalho aumenta?
    \end{itemize}

    \vspace{0.12cm}

    \begin{impactbox}{Hipótese}
        Para populações pequenas, o custo normalizado tenderá a ser estável.
        Para populações maiores, pode haver algumas degradaçãos associadas à hierarquia
        de memória.
    \end{impactbox}

\end{frame}

% ============================================================
% SLIDE 14
% ============================================================

\begin{frame}{Fatores experimentais}

    \vspace{0.08cm}

    \renewcommand{\arraystretch}{1.30}
    \begin{tabularx}{\textwidth}{>{\bfseries}c X X}
        Símbolo & Fator & Significado computacional \\
        \hline
        $N$ & Tamanho da população & Número de indivíduos processados por geração \\
        $L$ & Comprimento do cromossomo & Número de posições binárias por indivíduo \\
    \end{tabularx}

    \vspace{0.28cm}

    \[
        \text{carga nominal por execução} \approx N \times L \times G
    \]

    \vspace{0.10cm}

    \begin{questionbox}{Controle experimental}
        Taxa de mutação, probabilidade de crossover, método de seleção
        e número de gerações serão mantidos constantes.
    \end{questionbox}

\end{frame}

% ============================================================
% SLIDE 15
% ============================================================

\begin{frame}{Técnica de avaliação}

    \textbf{Técnica escolhida: medição em sistema real.}

    \vspace{0.08cm}

    \begin{columns}[T,onlytextwidth]

        \begin{column}{0.48\textwidth}
            \textbf{Ambiente experimental}

            \small
            \begin{itemize}
                \item Ubuntu 24.04.4 LTS / kernel 7.0.0-30-generic
                \item Intel Core i7-1355U (2 P-cores + 8 E-cores)
                \item RAM: 15 GiB
                \item Caches: L1 privada por núcleo; L2 privada/compartilhada
          conforme o tipo de núcleo; L3 12 MiB compartilhada
                \item C++ com GCC 13.3.0
                \item Compilação: otimização \texttt{-O3}
                \item Execução em uma única thread, fixada em um núcleo
            \end{itemize}
        \end{column}

        \begin{column}{0.48\textwidth}
            \textbf{Procedimento}

            \small
            \begin{itemize}
                \item Repetir cada configuração $(N,L)$.
                \item Registrar sementes pseudoaleatórias.
                \item Manter hardware, software e flags fixos.
            \end{itemize}

            \vspace{0.25cm}

            \begin{impactbox}{Ferramentas previstas}
                \small
                \begin{itemize}
                    \item Relógio monotônico: tempos.
                    \item \texttt{perf stat}: contadores de hardware.
                    \item \texttt{/usr/bin/time -v}: pico de memória RAM.
                \end{itemize}
            \end{impactbox}
        \end{column}

    \end{columns}

\end{frame}

% ============================================================
% SLIDE 16
% ============================================================

\begin{frame}{Métricas primárias}

    \begin{columns}[T,onlytextwidth]

        \begin{column}{0.50\textwidth}

            \textbf{Tempo por geração}

            \[
                T_g = \frac{T}{G}
            \]

            \begin{itemize}
                \item Unidade: ms/geração
                \item Lower is better
            \end{itemize}

        \end{column}

        \begin{column}{0.47\textwidth}

            \textbf{Vazão normalizada}

            \[
                Q = \frac{NLG}{T}
            \]

            \begin{itemize}
                \item Posições de cromossomo/s
                \item Higher is better
            \end{itemize}

        \end{column}

    \end{columns}

    \vspace{0.22cm}

    \begin{questionbox}{Interpretação}
        Se $Q$ cair conforme $N$ e $L$ aumentam, cada unidade nominal
        de trabalho está ficando mais cara.
    \end{questionbox}

\end{frame}

% ============================================================
% SLIDE 17
% ============================================================

\begin{frame}{Métricas explicativas}

    \begin{columns}[T,onlytextwidth]

        \begin{column}{0.48\textwidth}

            \textbf{Contadores de hardware}

            \begin{itemize}
                \item Ciclos de CPU
                \item Instruções executadas
                \item IPC
                \item Referências e faltas de cache
                \item Pico de memória RAM
            \end{itemize}

            \[
                IPC =
                \frac{\text{instruções}}
                     {\text{ciclos}}
            \]

        \end{column}

        \begin{column}{0.48\textwidth}

            \textbf{Decomposição do tempo}

            \begin{itemize}
                \item Avaliação da aptidão
                \item Seleção
                \item Crossover
                \item Mutação
            \end{itemize}

            \[
                P_{\text{fase}}
                =
                \frac{T_{\text{fase}}}
                     {T_{\text{total}}}
                \times 100\%
            \]

            \begin{center}
                \small
                Percentual do tempo gasto em cada etapa.
            \end{center}

        \end{column}

    \end{columns}

    \vspace{0.08cm}

    \begin{questionbox}{Registro das medições}
        Uma linha por execução em arquivo CSV, contendo
        $N$, $L$, $G$, semente, número da repetição e todas
        as métricas coletadas.
    \end{questionbox}

\end{frame}

% ============================================================
% SLIDE 18
% ============================================================

\begin{frame}{Resultado esperado}

    \textbf{Nosso estudo não pretende simplesmente concluir que cargas maiores demoram mais.}

    \vspace{0.18cm}

    \begin{itemize}
        \item Identificar se há mudança de regime no desempenho.
        \item Medir se o custo por unidade de trabalho permanece estável.
        \item Relacionar degradações a eventos de hardware quando possível.
        \item Descobrir qual fase do AG domina o tempo de execução.
    \end{itemize}

    \vspace{0.22cm}

    \begin{impactbox}{Mensagem principal}
        O objetivo é explicar o comportamento computacional do AG,
        não avaliar sua eficácia como algoritmo de otimização, pois este não é o foco da disciplina.
    \end{impactbox}

\end{frame}

% ============================================================
% SLIDE 19
% ============================================================

\begin{frame}{Resumo da Etapa 1}

    \renewcommand{\arraystretch}{1.28}
    \begin{tabularx}{\textwidth}{>{\bfseries}c X}
        Objeto computacional & AG sequencial simples, binário, com OneMax, medido sob variação de $N$ e $L$. \\
        \hline
        Método de análise & Medição em sistema real Linux, com instrumentação interna e contadores de hardware. \\
        \hline
        Métricas & Tempo por geração, vazão normalizada, ciclos, instruções, IPC, cache misses, memória e tempo por fase. \\
    \end{tabularx}

    \vspace{0.30cm}

    \begin{questionbox}{Por que nossa escolha se justifica?}
        O AG fornece uma carga computacional parametrizável. Ao alterar $N$ e $L$,
        podemos estudar escala dos dados, custo de execução e gargalos da máquina.
    \end{questionbox}

\end{frame}

\end{document}